\begin{table}
\begin{tcolorbox}[tabularx={X|p{2cm}|X|},enhanced,fonttitle=\bfseries\large,fontupper=\normalsize\sffamily,
colback=yellow!10!white,colframe=red!50!black,colbacktitle=blue!30!white,
coltitle=black,title=Drawing Options for Layered Graph Files (Part 1)]
\hline
Command & Arguments & Purpose \\
\hline
\hline
\verb'-paperheight'  &   h   &  Set paperheight for the latex document. The given value is measures the paperheight in tens of an inch. \\
\hline
\verb'-paperwidth'  &   w   &  Set paperwidth for the latex document. The given value is measures the paperwidth in tens of an inch. \\
\hline
\verb'-xin'  &   $a$   &  Assume input $x$-coordinates are in the interval $[0,a]$. Default value: 10000. \\
\hline
\verb'-yin'  &   $a$   &  Assume input $y$-coordinates are in the interval $[0,a]$. Default value: 10000. \\
\hline
\verb'-xout'  &   $a$   &  Assume output $x$-coordinates are in the interval $[0,a]$. Default value: 1000000. \\
\hline
\verb'-yout'  &   $a$   &  Assume output $y$-coordinates are in the interval $[0,a]$. Default value: 1000000. \\
\hline
\verb'-spanning_tree'  &     &  Place nodes according to a spanning tree. Default value: off.\\
\hline
\verb'-circle'  &     &  Circle all nodes. Default value: on.\\
\hline
\verb'-corners'  &     &  Draw corners at the outside of the drawing. Default value: off.\\
\hline
\verb'-radius'  &  $r$   &  Use radius $r$ for drawing circles around nodes. \\
\hline
\verb'-embedded'  &     &  Create latex headers for stand-alone latex files. Default value: off. \\
\hline
\verb'-sideways'  &     &  Create latex figure sideways. Default value: off.\\
\hline
\verb'-label_edges'  &     &  Label the edges. This option is relevant for the drawing of Schreier trees. Default value: off. \\
\hline
\verb'-x_stretch'  &  $s$   &  Apply $x$-axis scaling by a factor of $s$. Default value: $s=1.0$. \\
\hline
\verb'-y_stretch'  &  $s$   &  Apply $y$-axis scaling by a factor of $s$. Default value: $s=1.0$. \\
\hline
\end{tcolorbox}
\caption{\label{tab:draw:options}Drawing Options for Layered Graph Files (Part 1)}
\end{table}
